% !TeX root = ../测度论笔记.tex
% 由主文档通过 \input 载入；请编译主文档。
\setcounter{chapter}{9} % 保留原书章号。
\chapter{$L^p$ 空间与积分不等式}
\label{chap:mt-lp}

\section{为什么要把几乎处处相等的函数视为同一个元素}
若直接用 $\int|f-g|$ 衡量两个函数的距离，零测集上不同的函数会有距离零。因此需要先规定等价关系 $f\sim g$ 当且仅当 $f=g$ 几乎处处，再把每个等价类作为一个空间元素。实际计算时仍用 $f$ 表示等价类，但必须记住：一般 $L^p$ 元素的单点取值不是等价类本身确定的数据。

\begin{definition}
对 $1\leq p<\infty$，\emph{勒贝格空间（Lebesgue space）}$L^p(\mu)$ 是满足 $\int|f|^p\,\mathrm d\mu<\infty$ 的实值可测函数按几乎处处相等取商得到的空间，定义
\[
 \|f\|_p=\left(\int|f|^p\,\mathrm d\mu\right)^{1/p}.
\]
$L^\infty(\mu)$ 是本质有界可测函数的等价类空间，其范数为
\[
 \|f\|_\infty=\inf\{M\geq0:|f|\leq M\text{ 几乎处处}\},
\]
称为\emph{本质上确界（essential supremum）}范数。
\end{definition}
$L^\infty$ 中的“无穷”表示一个端点范数，而不是将 $|f|^\infty$ 直接拿去积分。对区间上的勒贝格测度，$\mathbf1_{\{0\}}$ 的本质上确界为零，但其通常的上确界为一。

\section{Hölder 与 Minkowski 不等式}
\begin{theorem}[Hölder 不等式]
设 $1\leq p,q\leq\infty$，$1/p+1/q=1$，约定 $1/\infty=0$。若 $f\in L^p(\mu)$、$g\in L^q(\mu)$，则 $fg\in L^1(\mu)$，且
\begin{equation}
 \int|fg|\,\mathrm d\mu\leq\|f\|_p\|g\|_q.
 \label{eq:mt-holder}
\end{equation}
\end{theorem}
\begin{proof}
当 $p=1,q=\infty$ 时，直接由 $|g|\leq\|g\|_\infty$ 几乎处处积分；另一端点相同。若 $1<p,q<\infty$，对 $a,b\geq0$ 有 Young 不等式 $ab\leq a^p/p+b^q/q$。固定 $b$ 后，对 $a$ 求极小值可验证此式：$a^p/p-ab$ 在 $a=b^{1/(p-1)}$ 处取最小值 $-b^q/q$。在两范数非零时，代入 $a=|f|/\|f\|_p$、$b=|g|/\|g\|_q$ 并积分，得到右侧为 $1/p+1/q=1$。若某个范数为零，对应函数几乎处处为零，结论同样成立。
\end{proof}

\begin{theorem}[Minkowski 不等式]
对 $1\leq p\leq\infty$ 和 $f,g\in L^p(\mu)$，
\begin{equation}
 \|f+g\|_p\leq\|f\|_p+\|g\|_p.
 \label{eq:mt-minkowski}
\end{equation}
\end{theorem}
\begin{proof}
$p=1,\infty$ 由逐点三角不等式得到。对 $1<p<\infty$，先由
$|f+g|^p\leq2^{p-1}(|f|^p+|g|^p)$ 知 $f+g\in L^p$；该点态估计来自 $t^p$ 的凸性。记 $q=p/(p-1)$，使用 Hölder 不等式可得
\[
 \|f+g\|_p^p
 \leq\int(|f|+|g|)|f+g|^{p-1}
 \leq(\|f\|_p+\|g\|_p)\|f+g\|_p^{p-1}.
\]
范数非零时相除，范数为零时结论直接成立。
\end{proof}

\section{有限测度与计数测度下的不同包含关系}
若 $0<\mu(S)<\infty$ 且 $1\leq p<q\leq\infty$，则
\begin{equation}
 \|f\|_p\leq\mu(S)^{1/p-1/q}\|f\|_q,\qquad L^q(\mu)\subseteq L^p(\mu).
 \label{eq:mt-finite-embedding}
\end{equation}
当 $q<\infty$ 时，对 $|f|^p$ 与常数 $1$ 使用指数 $q/p$ 及其共轭指数的 Hölder 不等式，再开 $p$ 次方即可；$q=\infty$ 时直接积分。若 $\mu(S)=0$，所有等价类均为零，包含关系仍成立。

在 $\mathbb N$ 的计数测度下，$L^p$ 记为 $\ell^p$，方向却是 $\ell^p\subseteq\ell^q$，$p<q$。事实上，$\sup_n|a_n|\leq\|a\|_p$，从而当 $q<\infty$ 时，
$\sum_n|a_n|^q\leq\|a\|_p^{q-p}\sum_n|a_n|^p$。两种方向不同，是因为数列每个坐标都有固定测度一，而有限区间上的函数可以在任意小集合上形成高峰。

\section{完备性与简单函数的稠密性}
\begin{theorem}[$L^p$ 的完备性]
对任意测度空间及 $1\leq p\leq\infty$，$L^p(\mu)$ 是完备赋范线性空间，即每个范数 Cauchy 序列都在该范数下收敛。
\label{thm:mt-lp-complete}
\end{theorem}
\begin{proof}
先设 $p<\infty$。从 Cauchy 序列 $(f_n)$ 中选子列，使 $\|f_{n_{k+1}}-f_{n_k}\|_p\leq2^{-k}$。令
$h_m=\sum_{k=1}^{m}|f_{n_{k+1}}-f_{n_k}|$。Minkowski 不等式给出 $\|h_m\|_p\leq1$；单调收敛作用于 $h_m^p$，说明 $h=\sum_{k=1}^{\infty}|f_{n_{k+1}}-f_{n_k}|$ 属于 $L^p$，故几乎处处有限。

因此该子列几乎处处绝对地累积变化，定义极限 $f$，并在例外零测集上取零。由 $|f|\leq|f_{n_1}|+h$ 可知 $f\in L^p$。对尾部级数重复同一估计，得到
\[
 \|f-f_{n_m}\|_p\leq\sum_{k=m}^{\infty}2^{-k}\longrightarrow0.
\]
原序列的 Cauchy 性再把子列收敛推广为整个序列收敛。

若 $p=\infty$，同样选取差的范数不超过 $2^{-k}$ 的子列，并合并每个差的点态界失效的零测集。在其补集上，差的级数被数值级数 $\sum_k2^{-k}$ 一致控制，故子列一致收敛于某个本质有界可测函数，尾部上界给出 $L^\infty$ 收敛。最后仍由 Cauchy 性处理整个序列。
\end{proof}

对 $p<\infty$，可积的简单函数在 $L^p$ 中稠密。对实值 $f$，将第\ref{chap:mt-functions}章的非负简单逼近用于 $|f|$，再乘以 $\operatorname{sgn}(f)$，其中符号函数在正值处取 $1$、负值处取 $-1$、零处取 $0$，得到简单函数 $s_n\to f$，且 $|s_n|\leq|f|$。它们与 $f$ 同号，故 $|f-s_n|^p\leq|f|^p$，控制收敛给出 $\|f-s_n\|_p\to0$。

$L^2$ 还具有内积 $\langle f,g\rangle=\int fg\,\mathrm d\mu$，Hölder 不等式保证该积分存在；完备性使它成为 Hilbert 空间。至此即可与项目中的泛函分析课程衔接，继续学习正交投影、连续线性泛函和算子。

\section{范数的几何含义与等号条件}
$L^1$ 范数累计绝对误差，$L^2$ 范数累计平方误差后开方，$L^\infty$ 范数控制几乎处处的最大误差。它们对尖峰的敏感程度不同。对高度为 $H$、支撑测度为 $a$ 的示性型函数，$\|H\mathbf1_A\|_p=H a^{1/p}$，而本质上确界为 $H$（假设 $a>0$）。因此，当支撑缩小时，有限阶范数可能变小，本质上确界却不变。

选择范数也就是选择允许的误差类型。平均预测误差适合积分范数，要求所有位置统一控制的任务更接近本质上确界；在偏微分方程中，能量通常对应平方积分，而函数本身与导数的积分控制又会组合成新的空间。不能仅因两个范数都能衡量大小，就认为它们的收敛概念或完备空间相同。

对 $1<p<\infty$，Hölder 不等式的等号条件可以从 Young 不等式逐点读出。若两个范数都非零，则等号成立当且仅当
\[
 \frac{|f|^p}{\|f\|_p^p}=\frac{|g|^q}{\|g\|_q^q}
 \quad\text{几乎处处}.
\]
因为归一化后的 Young 不等式之差非负，积分为零意味着它几乎处处为零，而数值 Young 等号恰要求 $a^p=b^q$。这是绝对值乘积积分的等号条件；若讨论 $|\int fg|$ 的等号，还必须保证乘积的符号不会发生抵消。

\section{连续紧支撑函数的稠密性}
\label{sec:mt-continuous-density}
简单函数已经在 $L^p$ 中稠密，但在欧氏空间上还希望用连续函数，甚至光滑函数逼近，以便应用微积分。这里证明实轴上连续紧支撑函数的稠密性，记这类函数组成的空间为 $C_c(\mathbb R)$。下标表示支撑紧，即函数在某个有界闭集之外为零。

\begin{theorem}
对 $1\leq p<\infty$，$C_c(\mathbb R)$ 在 $L^p(\mathbb R,\lambda)$ 中稠密。
\label{thm:mt-cc-density}
\end{theorem}
\begin{proof}
先考虑有限测度集合的示性函数。通过截断到大区间，可把问题归结为有界可测集 $E$。由正则性，选紧集 $K\subseteq E$ 与有界开集 $U\supseteq E$，使 $\lambda(U\setminus K)$ 任意小。若 $K$ 为空且 $E$ 测度已经足够小，可以直接用零函数逼近；否则定义
\[
 \varphi(x)=\frac{d(x,U^c)}{d(x,U^c)+d(x,K)},\qquad
 d(x,A)=\inf_{a\in A}|x-a|.
\]
紧集 $K$ 包含于开集 $U$，两者与补集之间的正分离保证分母处处非零。距离函数连续，因此 $\varphi$ 连续，且 $0\leq\varphi\leq1$，在 $K$ 上为一、在 $U$ 外为零，支撑包含于有界闭集 $\overline U$。于是
$\|\varphi-\mathbf1_E\|_p^p\leq\lambda(U\setminus K)$。

有限测度集合的有限线性组合可以逐项逼近；再用前面已证明的简单函数稠密性和三角不等式，便得到一般 $L^p$ 函数的结论。
\end{proof}
这一定理是测度正则性的一项分析后果。证明中先用紧集与开集控制集合误差，再用距离函数在二者之间连续过渡，最后以积分范数吸收过渡区域。近似函数不必在每个点都逼近原函数；允许把不连续性集中到总测度很小的区域里。

结论一般不能将 $p<\infty$ 改为 $p=\infty$。例如阶跃函数 $\mathbf1_{(0,\infty)}$ 不可能被连续函数在本质上确界范数下以小于 $1/2$ 的误差逼近：若误差小于某个 $c<1/2$，连续性将迫使近零点左侧的值不超过 $c$、右侧的值不小于 $1-c$，在零点取极限产生矛盾。小测度的过渡区域仍会被本质上确界检测到。

\section{可分性、无限测度空间与范数选择}
称度量空间\emph{可分（separable）}，是指存在可数稠密子集。对 $1\leq p<\infty$，$L^p(\mathbb R,\lambda)$ 可分：连续紧支撑函数可由有限个有理端点区间上的有理值阶梯函数在 $L^p$ 中逼近，而这样的阶梯函数只有可数多个。证明时先固定包含支撑的有界区间，利用一致连续性选取细网格，再将函数值近似为有理数；区间长度有限把一致误差转成积分误差。

相对地，$L^\infty(0,1)$ 不可分。集合 $\{\mathbf1_{(0,t)}:0<t<1\}$ 中任意两个不同元素的距离都是一。若存在可数稠密子集，每个这样的函数都应在该子集中某个元素的距离 $1/3$ 内；而同一个中心不能同时接近两个距离为一的函数，因此需要不可数多个中心，矛盾。这一差异会影响后续关于逼近、对偶空间和弱拓扑的研究。

在一般无限测度空间上，不同 $L^p$ 之间通常没有统一的包含方向。对实轴上的 $1\leq p<q<\infty$，取 $1/q<\alpha<1/p$，函数 $x^{-\alpha}\mathbf1_{(0,1)}$ 属于 $L^p$ 而不属于 $L^q$；函数 $(1+x)^{-\alpha}\mathbf1_{(1,\infty)}$ 则属于 $L^q$ 而不属于 $L^p$。前者的障碍位于局部奇点，后者的障碍位于远端尾部。有限测度空间上的包含关系，不能脱离背景测度直接使用。

\section{平移连续性与光滑逼近}
在实轴的勒贝格空间中，平移是保持范数的运算。记 $(\tau_hf)(x)=f(x-h)$，由平移不变性可得 $\|\tau_hf\|_p=\|f\|_p$。对 $1\leq p<\infty$，还有
$\|\tau_hf-f\|_p\to0$ 当 $h\to0$。这里的连续性是关于平移参数和积分范数而言，并不声称原函数本身逐点连续。

证明先对 $g\in C_c(\mathbb R)$ 进行：当 $|h|\leq1$，所有平移的支撑位于同一个有界区间中，一致连续性保证函数差的一致范数趋零，再乘以该区间长度的 $1/p$ 次方，就控制了 $L^p$ 范数。对一般 $f$，选连续紧支撑 $g$ 使 $\|f-g\|_p$ 小，再用
$\|\tau_hf-f\|_p\leq2\|f-g\|_p+\|\tau_hg-g\|_p$。
先固定逼近误差，再缩小平移量，得到结论。

取非负光滑紧支撑函数 $\rho$，支撑包含于 $[-1,1]$，且 $\int\rho=1$。例如在 $|x|<1$ 上取常数倍的 $\exp(-1/(1-x^2))$，在其余位置取零，再用归一化常数使积分为一。定义 $\rho_\varepsilon(x)=\varepsilon^{-1}\rho(x/\varepsilon)$，则其质量仍为一，但支撑缩到原点附近。这类函数称为\emph{磨光核（mollifier）}。

对 $f\in L^p(\mathbb R)$，卷积 $f*\rho_\varepsilon$ 将每个位置附近的函数值作局部平均。利用以 $\rho_\varepsilon(y)\,\mathrm dy$ 为概率测度的 Jensen 不等式，再用 Tonelli 定理，
\[
 \|f*\rho_\varepsilon-f\|_p^p
 \leq\int\rho_\varepsilon(y)\|\tau_yf-f\|_p^p\,\mathrm dy
 \leq\sup_{|y|\leq\varepsilon}\|\tau_yf-f\|_p^p\longrightarrow0.
\]
该估计适用于有限 $p\geq1$；内层积分的有限性由同一个非负积分估计在几乎处处意义下保证。

卷积后的函数光滑，是因为可以把对 $x$ 的导数移到核函数上。对局部的 $x$ 范围，各阶核导数统一有界且只涉及一个共同紧集上的 $f$，而 $L^p$ 函数在有限区间上属于 $L^1$，所以参数积分的控制条件成立。若先把 $f$ 截断到有界区间，再作磨光，就得到光滑紧支撑函数，从而证明 $C_c^\infty(\mathbb R)$ 在有限阶 $L^p$ 中稠密。

这一构造在偏微分方程中很重要：先对光滑近似进行微分和积分运算，再利用范数收敛将结论传回较粗糙的原函数。但近似的导数是否也收敛，需要额外的导数控制；仅有函数的 $L^p$ 收敛不能自动推出其导数收敛。选择正确的函数空间，正是为了把这些不同层次的控制同时记录下来。

\section{$L^2$ 中的正交与最小误差}
对实值 $L^2$ 函数，称 $f,g$ 正交，若 $\int fg\,\mathrm d\mu=0$。此时展开平方可得
$\|f+g\|_2^2=\|f\|_2^2+\|g\|_2^2$，这就是函数空间中的勾股关系。它不表示两个函数的支撑必须不交；支撑不交是充分条件，但正负抵消也可能使内积为零。

例如在 $[-1,1]$ 上，常数函数一与函数 $x$ 正交，因为后者是奇函数；两者却在几乎整个区间同时非零。若要用常数 $c$ 逼近平方可积函数 $f$，展开
$\int(f-c)^2=\int(f-\overline f)^2+\mu(S)(c-\overline f)^2$，
其中 $0<\mu(S)<\infty$、$\overline f=\mu(S)^{-1}\int f$，可见平均值是唯一的最佳常数近似。后面的条件期望会把“常数”替换成“关于指定信息可测的函数”，但最小平方误差的几何结构保持相同。

\section{习题}
\begin{enumerate}
 \item 对 $f(x)=x^{-\alpha}$，确定其在 $(0,1)$ 上属于 $L^p$ 的条件，其中 $\alpha>0$、$1\leq p<\infty$。
 \item 给出属于 $\ell^2$ 但不属于 $\ell^1$ 的数列，以及属于 $L^1(0,1)$ 但不属于 $L^2(0,1)$ 的函数。
 \item 证明 $\|f\|_p=0$ 当且仅当 $f=0$ 几乎处处，并解释取等价类在范数定义中的作用。
\end{enumerate}

